\chapter*{\Large \center Abstract}

% GUIDANCE (about 250 words; write it LAST, after the chapters):
%   1. One sentence of context: at Festival of the Mind 2024 visitors typed how they felt and a
%      text-to-image model drew each entry as a monochrome line drawing (\nTotal records,
%      \nAnalysed analysed).
%   2. The question: does the emotion a reader recovers from the drawing agree with the emotion a
%      reader recovers from the words — measured as SIMILARITY on a ladder of four measures, not as
%      an exact match, and always against a random-pairing null.
%   3. Method in two sentences: prompt reader (GoEmotions, 28 -> 7 families); image side described
%      by BLIP-VQA with five factual questions and read by zero-shot NLI; four measures; null and z;
%      clustered bootstrap CIs and clustered permutation tests because \nRepeats records repeat a prompt.
%   4. Headline numbers (all macros): emotional prompts valence-sign \FullEmoSignObs\ vs null
%      \FullEmoSignNull\ ($z$ \FullEmoSignZ); face-classified \FullFaceSignObs\ vs scene-classified
%      \FullSceneSignObs, difference \DiffSignPp\ (95\% CI \DiffSignCi, $p$ = \DiffSignP);
%      neutral-labelled prompts at their null (VA \FullNeuVaObs\ vs \FullNeuVaNull); describe-then-
%      extract vs direct CLIP like for like, family \McDeFam\ vs \McClipFam\ (\McDiffFamPp);
%      repeated prompts \ConsFam\ vs null \ConsFamNull\ ($z$ \ConsFamZ).
%   5. The honest turn: the split is the describer's, not a human's; the describer answers an
%      expression for \DescNoPeopleExpr\ drawings in which it reports no people, and where its answer
%      names an emotion the image reader copies it \DescCopyStrict\ of the time — so the finding is
%      about a measurement chain, and a human-rated study is the next step.
%   6. Close with what is contributed: the corpus-plus-question, the null-referenced similarity
%      ladder, the audits, and the open demonstrator.
